\section{Related Work}
\label{sec:related}

\paragraph*{Supply-chain incident analysis.}
Detailed post-mortems of individual registry compromises establish the attack
patterns we set out to
measure~\citep{ferrante2021incidents,mbeki2022typosquat}. They are necessarily
retrospective and selected on having been detected, so they cannot speak to
prevalence, which is the gap our longitudinal crawl addresses.

\paragraph*{Static analysis of package code.}
Several systems scan package archives for suspicious constructs before
publication~\citep{ohara2020staticscan}. They are fast enough to run on every
upload and they catch unobfuscated payloads. Every sample in our malicious set
would evade the published feature sets, because each used at least runtime
string assembly, and Listing~\ref{lst:hook} shows how little that takes.

\paragraph*{Dynamic analysis and sandboxing.}
Sandboxed execution of untrusted package code has been proposed for runtime
imports~\citep{vanRhijn2019sandbox} and for build
steps~\citep{szymanska2022buildsandbox}. Our contribution is to apply it at
registry scale over a long window and to report what the population looks like,
rather than to propose the sandbox as a novel mechanism.

\paragraph*{Reproducible and hermetic builds.}
Hermetic build systems eliminate install-time hooks by
construction~\citep{adeyemi2021hermetic}, and where an organisation can adopt
one, the problem described here does not arise. Adoption in the ecosystems we
study is low and the transition cost is high, so we consider hermetic builds
the right long-term answer and sandboxed hooks the deployable one.

\paragraph*{Account security signals.}
Work on detecting account takeover from behavioural
telemetry~\citep{quintero2020takeover} supplies the mechanism behind our
strongest result. We contribute the observation that the specific conjunction
of a recovery-email change with a hook-bearing publication is both highly
predictive and already available to every registry.
